% ===== PRÉAMBULE CANONIQUE — à coller VERBATIM en tête de CHAQUE .tex =====
\documentclass[11pt,a4paper]{article}
\usepackage[utf8]{inputenc}\usepackage[T1]{fontenc}\usepackage{lmodern}
\usepackage[french]{babel}
\usepackage{amsmath,amssymb,mathtools}\usepackage{icomma}
\usepackage[version=4]{mhchem}          % \ce{...} : équations & formules
\usepackage{siunitx}\sisetup{locale=FR} % \qty{}{}, \si{}, \ang{}
\usepackage{chemfig}                    % \chemfig{...} : structures développées
\usepackage{xcolor}\usepackage{colortbl}
\usepackage{tikz}\usepackage{pgfplots}\pgfplotsset{compat=1.18}
\usepackage[most]{tcolorbox}\usepackage{enumitem}\usepackage{array,booktabs}
\usepackage[a4paper,margin=2.2cm]{geometry}
\usetikzlibrary{arrows.meta,calc,angles,quotes,positioning,patterns,backgrounds,shapes.geometric}
\definecolor{primary}{RGB}{222,49,99}\definecolor{primarydark}{RGB}{150,22,55}
\definecolor{defcol}{RGB}{0,114,178}\definecolor{methcol}{RGB}{52,92,148}
\definecolor{excol}{RGB}{0,135,90}\definecolor{attncol}{RGB}{200,40,55}
\tcbset{boxbase/.style={enhanced,breakable,boxrule=0.9pt,arc=2.5pt,
  left=3mm,right=3mm,top=2.3mm,bottom=2.3mm,fonttitle=\bfseries,coltitle=white}}
% --- boîtes TUTORIEL (noms EXACTS, ne pas renommer) ---
\newtcolorbox{cle}[1][]{boxbase,colback=defcol!6,colframe=defcol,
  title={\textsc{À retenir}\ifx\relax#1\relax\relax\else\textmd{ : \itshape #1}\fi}}
\newtcolorbox{methode}[1][]{boxbase,colback=methcol!7,colframe=methcol,sharp corners,
  title={\textsc{Méthode}\ifx\relax#1\relax\relax\else\textmd{ --- \itshape #1}\fi}}
\newtcolorbox{exemplebox}[1][]{boxbase,colback=excol!5,colframe=excol,
  title={\textsc{Exemple}\ifx\relax#1\relax\relax\else\textmd{ : \itshape #1}\fi}}
\newtcolorbox{piege}[1][]{boxbase,colback=attncol!6,colframe=attncol,
  title={\textsc{Piège}\ifx\relax#1\relax\relax\else\textmd{ : \itshape #1}\fi}}
\newtcolorbox{remarque}[1][]{boxbase,colback=black!4,colframe=black!45,
  title={\textsc{Remarque}\ifx\relax#1\relax\relax\else\textmd{ : \itshape #1}\fi}}
% --- boîtes EXERCICE / CORRIGÉ (noms EXACTS, auto-comptées) ---
\newtcolorbox[auto counter]{exo}[1][]{boxbase,colback=primary!4,colframe=primary,
  title={\textsc{Exercice}~\thetcbcounter\ifx\relax#1\relax\relax\else\textmd{ --- \itshape #1}\fi}}
\newtcolorbox[auto counter]{corrige}[1][]{boxbase,colback=excol!4,colframe=excol,
  title={\textsc{Corrigé}~\thetcbcounter\ifx\relax#1\relax\relax\else\textmd{ --- \itshape #1}\fi}}
\newcommand{\R}{\mathbb{R}}\newcommand{\N}{\mathbb{N}}\newcommand{\Z}{\mathbb{Z}}\newcommand{\ds}{\displaystyle}
% (un \usepackage{fancyhdr}+en-tête est possible ; il est ignoré côté web)
% =========================================================================
\begin{document}
% THEME_TITLE: Gaz nobles, règle de l'octet, familles et tendances périodiques

La position d'un élément « raconte » son comportement. Ce thème relie les \textbf{électrons de valence}
à la \textbf{règle de l'octet}, aux \textbf{familles} chimiques et aux grandes \textbf{tendances} du
tableau périodique.

\begin{cle}[gaz nobles et règle de l'octet]
Les \textbf{gaz nobles} (He, Ne, Ar\dots) ont une couche de valence \textbf{saturée} : $8$ électrons
(un \textbf{octet}) pour tous, sauf l'hélium qui se contente de $2$ (un \textbf{duet}). Cette saturation
les rend très \emph{stables}. \textbf{Règle de l'octet :} les autres atomes gagnent, perdent ou
partagent des électrons pour \emph{atteindre la configuration du gaz noble le plus proche}.
\end{cle}

\begin{cle}[familles et ion formé]
Une \textbf{famille} (colonne) regroupe les éléments de \emph{même nombre d'électrons de valence} :
\begin{itemize}[leftmargin=1.5em,topsep=2pt,itemsep=1pt]
  \item \textbf{alcalins} ($1$ e$^-$ de valence) : le perdent $\to$ cation $\ce{X^{+}}$ ;
  \item \textbf{alcalino-terreux} ($2$) : cation $\ce{X^{2+}}$ ;
  \item \textbf{halogènes} ($7$) : gagnent $1$ e$^-$ $\to$ anion $\ce{X^{-}}$ ;
  \item \textbf{gaz nobles} ($8$, ou $2$ pour He) : saturés, quasi inertes.
\end{itemize}
\end{cle}

\begin{methode}[prévoir l'ion d'un élément]
Compter les électrons de valence $v$ :
\begin{enumerate}[leftmargin=1.7em,topsep=2pt,itemsep=1pt]
  \item $v \le 3$ (métaux) : l'atome \emph{perd} ses $v$ électrons $\to$ cation de charge $+v$ ;
  \item $v \ge 5$ (non-métaux) : l'atome \emph{gagne} $8-v$ électrons $\to$ anion de charge $-(8-v)$ ;
  \item dans les deux cas, l'ion adopte la configuration du gaz noble le plus proche.
\end{enumerate}
\end{methode}

\begin{cle}[tendances périodiques]
\begin{center}
\begin{tikzpicture}[scale=0.9]
  \fill[primary!7] (0,0) rectangle (6,2.4);
  \draw[primary,line width=1pt] (0,0) rectangle (6,2.4);
  \draw[-{Stealth[length=2.4mm]},defcol,line width=1.1pt] (0.2,2.75) -- (5.8,2.75);
  \node[defcol,font=\scriptsize,above] at (3,2.72) {Électronégativité \& énergie d'ionisation $\nearrow$};
  \draw[-{Stealth[length=2.4mm]},defcol,line width=1.1pt] (-0.35,0.2) -- (-0.35,2.2);
  \draw[{Stealth[length=2.4mm]}-,attncol,line width=1.1pt] (0.2,-0.35) -- (5.8,-0.35);
  \node[attncol,font=\scriptsize,below] at (3,-0.32) {Rayon atomique $\nwarrow$ (augmente vers la gauche/le bas)};
  \node[font=\scriptsize,primarydark] at (0.95,1.2) {métaux};
  \node[font=\scriptsize,excol] at (5.05,1.2) {non-métaux};
\end{tikzpicture}
\end{center}
\vspace{-0.3em}
L'\textbf{électronégativité} et l'\textbf{énergie d'ionisation} augmentent vers la \textbf{droite} et le
\textbf{haut} ; le \textbf{rayon atomique} varie en sens inverse (vers la gauche et le bas). Le fluor
(coin haut-droite) est le plus électronégatif.
\end{cle}

\begin{piege}[l'hélium et la bonne charge]
L'hélium vise un \textbf{duet} ($2$), pas un octet. Et un atome forme l'ion qui atteint \emph{vraiment}
un gaz noble : $\ce{Ba^{2+}}$ (config.\ du xénon), pas $\ce{Ba^{+}}$ ; les métaux de transition (Cu, Fe\dots)
ne suivent pas simplement cette règle.
\end{piege}

\end{document}
